\documentclass[11pt]{article}

% ===========================================================================
% Preamble: packages, theorem environments, macros.
% Style mirrors the Zinc+ paper / other Zinc+ docs.
% ===========================================================================

\usepackage[a4paper,margin=1in]{geometry}
\usepackage[utf8]{inputenc}
\usepackage[T1]{fontenc}
\usepackage{lmodern}
\usepackage{microtype}

\usepackage{amsmath,amssymb,amsthm,mathtools}
\usepackage{booktabs}
\usepackage{array}
\usepackage{tabularx}
\usepackage{enumitem}
\usepackage{xcolor}
\usepackage{listings}
\usepackage[colorlinks=true,
            linkcolor={teal!70!black},
            citecolor={teal!70!black},
            urlcolor={teal!70!black}]{hyperref}
\usepackage{cleveref}

% --- Theorem environments --------------------------------------------------
\theoremstyle{plain}
\newtheorem{theorem}{Theorem}[section]
\newtheorem{lemma}[theorem]{Lemma}
\newtheorem{proposition}[theorem]{Proposition}
\newtheorem{corollary}[theorem]{Corollary}
\newtheorem{claim}[theorem]{Claim}

\theoremstyle{definition}
\newtheorem{definition}[theorem]{Definition}
\newtheorem{example}[theorem]{Example}

\theoremstyle{remark}
\newtheorem{remark}[theorem]{Remark}

% --- Operators, macros -----------------------------------------------------
\DeclareMathOperator{\GF}{GF}
\DeclareMathOperator{\Tr}{Tr}

\newcommand{\F}{\mathbb{F}}
\newcommand{\Q}{\mathbb{Q}}
\newcommand{\Z}{\mathbb{Z}}
\newcommand{\N}{\mathbb{N}}

% Cantor basis vectors
\newcommand{\bb}{\beta}

% Modulus polynomial: F_2[X]/f for the field, Z[X]/ftilde for the lift.
\newcommand{\flift}{\tilde f}

% Lifted ring shorthand: Z[X]/(ftilde)
\newcommand{\Rlift}{\Z[X]/\flift}

% Inline code style for short tokens
\newcommand{\code}[1]{\texttt{#1}}

% Math-mode shorthand for the codebase's `Int<N>` type.
\newcommand{\Int}{\mathtt{Int}}

% Define Rust as a listings language (the `listings` package doesn't ship
% one by default).
\lstdefinelanguage{Rust}{
  keywords={as, async, await, break, const, continue, crate, dyn, else, enum,
            extern, false, fn, for, if, impl, in, let, loop, match, mod, move,
            mut, pub, ref, return, self, Self, static, struct, super, trait,
            true, type, unsafe, use, where, while},
  morekeywords={i8, i16, i32, i64, i128, u8, u16, u32, u64, u128, usize, isize,
                bool, char, str, String, Vec, Option, Result, Box},
  sensitive=true,
  morecomment=[l]{//},
  morecomment=[s]{/*}{*/},
  morestring=[b]",
}

\lstset{
  basicstyle=\ttfamily\small,
  breaklines=true,
  showstringspaces=false,
  frame=single,
  framesep=4pt,
  xleftmargin=4pt,
  xrightmargin=4pt,
  keywordstyle=\color{teal!70!black}\bfseries,
  commentstyle=\color{gray}\itshape,
  stringstyle=\color{teal!50!black},
}

% Hyperref-friendly section refs
\renewcommand{\sectionautorefname}{Section}
\renewcommand{\subsectionautorefname}{Section}
\renewcommand{\subsubsectionautorefname}{Section}


\title{A Binary Additive Reed--Solomon Code for Zip$^+$, Structurally Lifted to $\Z[X]/\flift$ \\
       \large Design, implementation and integration notes}
\author{Albert Garreta \\[2pt] \normalsize Nethermind Research}
\date{\today}

\begin{document}

\maketitle

\begin{center}
\textcolor{blue}{\textit{AI-generated implementation report. Reviewed by author.}}
\end{center}
\vspace{0.5em}

\begin{abstract}
We design and implement a linear code for the Zip$^+$ PCS that uses
an additive Reed--Solomon FFT over $\GF(2^{16}) = \F_2[X]/f$ as its
underlying structure, but applies that algorithm \emph{structurally
in the lifted integer ring} $\Z[X]/\flift$, where $\flift$ is the
$\Z$-polynomial with the same coefficients as $f$. The construction
is the binary-field analogue of the IPRS code already shipped in
Zinc$^+$, which lifts a multiplicative Reed--Solomon NTT over
$\F_{65537}$ to a pseudo-NTT over $\Z$.

The encoder uses a \emph{radix-8} evaluator: each meta-stage applies
an $8 \times 8$ $\GF(2^{16})$ Vandermonde matrix whose entries are
precomputed in the binary field and lifted to $\{0,1\}$-coefficient
polynomials, with the matrix--vector product carried out in
$\Z[X]/\flift$. Computing the Vandermonde in $\GF(2^{16})$ rather than
composing lifted twiddles in $\Z$ is what keeps the lifted codeword
coefficients near $2^{26}$ ($\sim 28$ signed bits) instead of the
$\sim 2^{47}$ a radix-2 $\Z$-lift produces --- a $\sim 21$-bit
reduction that lets each codeword cell pack into $32$ bits on the wire.

The code is integrated end-to-end into Zip$^+$'s commit / prove /
verify: the verifier-side field is $F = \F_p[X]/\flift$ (a degree-$16$
extension of $\F_p$ when $\flift$ is irreducible mod $p$), and the
row-batching challenges are small $4$-bit-coefficient polynomials. We
describe the algorithm, the Rust implementation in
\code{zip-plus/src/code/binary\_add\_fft\_16/}, the protocol
integration, and the design decisions behind the shipped encoder.
\end{abstract}

\tableofcontents

\section{Overview and motivation}
\label{sec:overview}

Zip$^+$ is the polynomial commitment scheme used by Zinc$^+$. It commits
to a multilinear polynomial whose evaluations live in some semiring
$\mathrm{Eval}$. The prover encodes each row of the evaluation matrix
with a linear error-correcting code, arranges the encoded rows into
a matrix, and builds a Merkle tree over the columns. The two
linear-code families already shipped are:

\begin{itemize}
  \item \textbf{IPRS}: a radix-8 pseudo Reed--Solomon NTT over
    $\F_{65537}$, structurally lifted to $\Z$ so that the encoder
    operates on small integers and avoids field reductions on the
    prover hot path. This is the workhorse for the SHA / ECDSA
    integer-side lanes.
  \item \textbf{RAA}: a repeat-accumulate-accumulate code with two
    random permutations. Cheaper than IPRS but with weaker distance
    properties.
\end{itemize}

\paragraph{The motivation.}
SHA-256 and ECDSA witnesses are \emph{natively binary}: each evaluation
is most naturally represented as a polynomial of degree $<32$ over
$\F_2$ (\code{BinaryPoly<32>} in the codebase). When these witnesses
are passed through IPRS, the binary nature of the input is invisible
to the encoder; the algebra runs over $\F_{65537}$ even though every
input has only $\{0,1\}$ coefficients. A natural question is
whether a Reed--Solomon code defined directly over the binary
extension field $\GF(2^{32})$ would be a tighter fit.

\paragraph{The construction at a glance.}
Pick an irreducible $f \in \F_2[X]$ of degree $32$. Let $\flift$ be
the integer polynomial with the same monomials and coefficients in
$\{0,1\} \subset \Z$. Take a standard additive Reed--Solomon FFT
algorithm over $\GF(2^{32}) = \F_2[X]/f$ --- we use the LCH
(Lin--Chung--Han) variant with the Cantor basis, since the latter
makes the subspace polynomials $F_2$-linearised, hence
$\{0,1\}$-coefficient even after the lift. Now apply the
\emph{exact same algorithm}, with the same twiddles, but with every
arithmetic operation performed in $\Rlift$ instead of $\GF(2^{32})$.
Inputs (\code{BinaryPoly<32>}) lift trivially: each $\F_2$-coefficient
becomes a $\Z$-coefficient in $\{0,1\}$.

The resulting code has two pleasant properties:
\begin{enumerate}
  \item \textbf{Mod-2 commutativity.} Reducing the lifted codeword
    coefficient-wise modulo $2$ recovers exactly the
    $\GF(2^{32})$-RS codeword. This is the soundness witness we
    rely on to argue that the lifted code inherits good distance
    from the original RS code, modulo the lift.
  \item \textbf{Sparse-twiddle multiplication.} Every twiddle in the
    FFT is a Cantor-basis vector or a subspace-polynomial value at
    a basis vector --- all of these are lifted $\GF(2^{32})$ elements
    with $\{0,1\}$ coefficients. Multiplying any element of
    $\Rlift$ by such a twiddle is a sparse shift-and-add, no
    coefficient multiplications.
\end{enumerate}

\paragraph{Scope of this document.}
We cover (i) the algorithm, (ii) the Rust implementation in
\code{zip-plus/src/code/binary\_add\_fft/}, (iii) the
\code{ZipTypes} plumbing that makes the no-alpha-projection variant
of Zip$^+$ admissible, (iv) the protocol-level gap that currently
prevents end-to-end prove/verify, (v) the optimization journey from
naive to $\sim 26\times$-faster commit, and (vi) measured proof-size
properties (raw and zstd-3 compressed) of column openings.

\section{The algorithm}
\label{sec:algorithm}

\subsection{Field and ring choices}
\label{sec:field-choices}

Fix
\[
  f(X) \;=\; X^{16} + X^{5} + X^{3} + X^{2} + 1 \;\in\; \F_2[X],
\]
an irreducible polynomial of degree~$16$. Define
$\GF(2^{16}) = \F_2[X]/f$ and $\flift \in \Z[X]$ to be the integer
polynomial with the same set of monomials and the same $\{0,1\}$
coefficients. The lifted ring is
\[
  \Rlift \;=\; \Z[X]/\flift,
\]
a $16$-dimensional $\Z$-module on which addition is component-wise and
multiplication is schoolbook poly-multiplication followed by reduction
modulo $\flift$. The reduction is
\[
  X^{16} \;\equiv\; X^{5} + X^{3} + X^{2} + 1 \pmod{\flift},
\]
so any product whose temporary degree exceeds $15$ is folded back by
contributing to four low positions per high monomial.

\subsection{The Cantor basis}
\label{sec:cantor}

The Cantor basis of $\GF(2^{16})$ is the sequence
$\bb_0, \bb_1, \ldots, \bb_{15}$ defined recursively by
\begin{equation}
  \bb_0 = 1, \qquad \bb_i^2 + \bb_i = \bb_{i-1} \quad \text{for } i \ge 1.
  \label{eq:cantor-recurrence}
\end{equation}
The recurrence is solvable because $\Tr(\bb_{i-1}) = 0$ at every step
($\Tr(\bb_i^2 + \bb_i) = \Tr(\bb_i)^2 + \Tr(\bb_i) = 0$ in $\F_2$).
For~$i = 0$, $\Tr(1) = 16 \cdot 1 = 0$ since the extension degree is
even.

The two crucial properties for the additive FFT are:
\begin{itemize}
  \item $V_i := \mathrm{span}_{\F_2}(\bb_0, \ldots, \bb_{i-1})$ is a
    chain of $\F_2$-subspaces of $\GF(2^{16})$, with
    $\dim V_i = i$.
  \item The associated subspace polynomial
    \[
      s_i(X) \;=\; \prod_{\alpha \in V_i} (X + \alpha)
    \]
    is $\F_2$-linearised: only monomials of degree $2^j$ for
    $j = 0, \ldots, i$ appear, all coefficients lie in $\F_2$, and
    the leading coefficient is~$1$. The Cantor basis additionally
    guarantees $s_i(\bb_i) = 1$, which simplifies the
    recurrence to
    \begin{equation}
      s_{i+1}(X) \;=\; s_i(X)^2 + s_i(X).
      \label{eq:cantor-subspace-recurrence}
    \end{equation}
\end{itemize}

Crucially for the structural lift: every $s_i$ coefficient is in
$\F_2$, so its lift to $\Z$ is in $\{0, 1\}$. Combined with the fact
that Cantor basis vectors themselves are elements of $\GF(2^{16})$,
every twiddle the FFT uses --- and hence every entry of the radix-8
Vandermonde matrices below --- is a $\{0, 1\}$-coefficient polynomial
once lifted to $\Rlift$.

\subsection{The novel polynomial basis}
\label{sec:lch}

Let $m = \log_2 N$, where $N$ is a power of $2$ at most $2^{16}$.
Express the input polynomial $P(X)$ of degree~$< N$ in the
\emph{novel polynomial basis}
\[
  X_i(X) \;=\; \prod_{j \in \{0..m-1\}\,:\, i_j = 1} s_j(X),
\]
where $i_j$ is the $j$-th bit of $i$. Then $P = \sum_{i=0}^{N-1} c_i \, X_i$,
and the forward FFT maps the coefficient vector
$(c_0, \ldots, c_{N-1})$ to the evaluation vector $(P(v_0), \ldots,
P(v_{N-1}))$ at the points $v_k = \sum_l k_l \bb_l \in V_m$. Data is
in natural order at both input and output --- no bit-reversal
permutation is required.

The textbook radix-2 LCH FFT computes this in $m$ stages: at stage $i$,
each block of $2^{i+1}$ entries is processed by $2^i$ butterflies, each
mixing a pair $(a,b)$ via a single twiddle $t = s_i(\omega)$:
$a' = a + t\cdot b$, $b' = a' + b$. We do \emph{not} use the radix-2
form directly --- see below.

\subsection{The radix-8 evaluator}
\label{sec:radix8}

The shipped encoder uses a \textbf{radix-8} variant. Three consecutive
radix-2 stages are merged into one \emph{meta-stage}: instead of
$m$ stages of $2\times 2$ butterflies, the FFT runs $m/3$ meta-stages,
each applying an $8 \times 8$ Vandermonde matrix per block of $8$
codeword positions. (This requires $3 \mid m$; the linear-code
constructor rejects parameters that violate it.)

Concretely, meta-stage $K$ covers the three subspace polynomials
$s_{3K}, s_{3K+1}, s_{3K+2}$. For each block $B$ of $8$ positions
the $8 \times 8$ matrix $V_{K,B}$ has entries
\[
  V_{K,B}[j][i] \;=\; X^{(K)}_i\!\left( v_{K,B,j} \right)
  \;=\;
  \prod_{a \,:\, i_a = 1} s_{3K+a}\!\left( v_{K,B,j} \right),
\]
the $i$-th novel-basis polynomial of the $3$-dimensional
sub-recursion, evaluated at the $j$-th point of that block. The
meta-stage output is the per-block matrix--vector product
\[
  \mathrm{out}[j] \;=\; \sum_{i=0}^{7} V_{K,B}[j][i] \cdot \mathrm{in}[i].
\]
Composed over the $m/3$ meta-stages, this evaluates exactly the same
multilinear extension as the radix-2 FFT --- over $\GF(2^{16})$ the two
agree identically.

\paragraph{Why radix-8, and why it matters for the lift.}
The radix-8 form is the load-bearing design choice for this code,
\emph{not} a peephole optimization. The reason is the magnitude of the
lifted codeword coefficients (\autoref{sec:optimization}):

\begin{itemize}
  \item In a radix-2 $\Z$-lift, the codeword is computed as a chain of
    $m$ sequential multiplications in $\Rlift$ by lifted twiddles. Each
    multiplication is a $\Z$-arithmetic \code{mul + reduce-mod-}$\flift$,
    so $\Z$-carries accumulate across all $m$ stages. Empirically the
    coefficients reach $\sim 2^{47}$ at $\code{codeword\_len} = 4096$.
  \item In the radix-8 form, each $V_{K,B}$ is computed \emph{in
    $\GF(2^{16})$} --- where the product of three subspace-polynomial
    values is just another $\GF(2^{16})$ element --- and only then
    lifted, so every Vandermonde entry has $\{0,1\}$ coefficients. The
    FFT then composes $m/3$ such $\{0,1\}$-coefficient matvecs in
    $\Rlift$. With only $m/3$ carry-accumulating meta-stages instead of
    $m$, and a $\{0,1\}$-coefficient multiplier at each, the codeword
    coefficients stay near $2^{26}$ ($\sim 28$ signed bits) --- a
    $\sim 21$-bit reduction over the radix-2 $\Z$-lift.
\end{itemize}

This is a genuine change of code: the radix-8-with-GF-lifted-Vandermonde
encoder produces a \emph{different} lifted code than a radix-2 $\Z$-lift
would. The two agree modulo $2$ (both reduce to the standard
$\GF(2^{16})$ RS codeword, see below), but differ in the higher
$\Z$-bits, because precomputing $V_{K,B}$ in $\GF(2^{16})$ discards the
inter-stage $\Z$-carries that a sequential radix-2 lift would keep.
Choosing the radix-8 code is choosing the variant with the small
coefficients.

\subsection{The structural lift and mod-2 commutativity}
\label{sec:lift}

The lifted FFT runs the radix-8 algorithm with arithmetic in $\Rlift$
rather than $\GF(2^{16})$. The Cantor basis, subspace polynomials, and
the $8 \times 8$ Vandermonde matrices are all computed once in
$\GF(2^{16})$ and then \emph{lifted} coefficient-by-coefficient from
$\F_2$ to $\{0, 1\} \subset \Z$. The inputs are \code{BinaryPoly<16>}
elements, whose $\F_2$-coefficients also lift trivially.

Each per-block matvec is a sequence of $\Rlift$-additions and
multiplications by lifted, $\{0,1\}$-coefficient Vandermonde entries.
The multiplication is a \emph{sparse shift-and-accumulate}: for every
set bit of the entry, add the operand shifted by that many positions
to a length-$31$ unreduced accumulator, then fold the accumulator back
to length~$16$ via the $\flift$ reduction.

The soundness witness is the commutative diagram: let
$\pi_2 : \Rlift \to \GF(2^{16})$ be coefficient-wise mod-$2$ reduction.
Then for any input vector $\mathbf c$,
\[
  \pi_2\!\bigl( \mathrm{LiftedRadix8FFT}(\mathrm{lift}(\mathbf c)) \bigr)
  \;=\;
  \mathrm{FFT}_{\GF(2^{16})}(\mathbf c).
\]
This holds because mod-$2$ reduction is a ring homomorphism that
commutes with addition and the $\flift$ reduction (as
$f \equiv \flift \pmod 2$), the lifted Vandermonde entries reduce mod
$2$ to the $\GF(2^{16})$ entries they were lifted from, and the
$\GF(2^{16})$ radix-8 evaluator computes exactly the standard
$\GF(2^{16})$ additive-FFT codeword. So the lifted code reduces mod $2$
to a genuine $\GF(2^{16})$ Reed--Solomon code, which is the basis for
inheriting that code's distance.

\section{Implementation}
\label{sec:implementation}

\subsection{Module layout}

The implementation lives in
\code{zip-plus/src/code/binary\_add\_fft\_16/} and mirrors the
existing IPRS layout in \code{zip-plus/src/code/iprs/}:

\begin{itemize}
  \item \code{basis.rs}: a $\GF(2^{16})$ \code{Gf2\_16} type (a
    \code{u16} newtype) with addition (XOR), multiplication
    (carry-less multiply + reduce mod $f$), the half-trace solver
    (Gaussian elimination over $\F_2$) for $Y^2 + Y = \beta$, and
    Cantor basis generation.
  \item \code{params.rs}: a \code{Config16} trait declaring the
    extension degree and irreducible $f$, plus the precomputed
    Cantor basis, evaluation points, and subspace polynomials
    (stored sparsely as coefficient lists at the $\{2^j\}$
    positions).
  \item \code{ring\_ops.rs}: the lifted-ring arithmetic. Exports
    \code{reduce\_mod\_ftilde\_16}, the lifted-twiddle
    multiplication primitive, and the \code{FftRingElement16} trait.
  \item \code{radix8.rs}: the radix-8 evaluator ---
    \code{Radix8FftParams16}, which precomputes the per-meta-stage
    $8 \times 8$ Vandermonde tables, and \code{additive\_fft\_radix8\_16},
    the kernel.
  \item \code{ext\_field.rs}: the verifier-side field
    $F = \F_p[X]/\flift$ (\code{Gf2\_16Ext<P>}) and the
    $4$-bit-coefficient polynomial challenge type \code{Bit4Poly16}.
    See \autoref{sec:protocol-integration}.
  \item \code{packed\_cw.rs}: \code{PackedI64Poly16<BITS>}, a
    codeword/combined-row cell type with a bit-packed wire encoding
    (\autoref{sec:packed-cw}).
\end{itemize}

The top-level \code{binary\_add\_fft\_16.rs} exposes the public
\code{BinaryAddFft16Code} struct and its \code{LinearCode<Zt>} impl.

\subsection{The {\tt FftRingElement16} trait}
\label{sec:fft-ring-element}

The kernel is polymorphic over a \code{FftRingElement16} trait that
exposes the primitives the radix-8 matvec needs:

\begin{lstlisting}[language=Rust]
pub trait FftRingElement16: Sized + Clone + Send + Sync {
    fn fft_zero() -> Self;
    fn fft_add(&self, other: &Self) -> Self;
    fn fft_mul_lifted_gf16(&self, twiddle: Gf2_16) -> Self;

    // Batched-reduction path for the 8x8 matvec: accumulate the
    // eight lifted-twiddle products in an UNREDUCED buffer, then
    // reduce mod ftilde exactly once per output cell.
    type UnreducedAcc: Send + Sync + Clone;
    fn fft_unreduced_zero() -> Self::UnreducedAcc;
    fn fft_acc_mul_lifted_gf16(
        input: &Self, twiddle: Gf2_16, acc: &mut Self::UnreducedAcc,
    );
    fn fft_finalize_acc(acc: Self::UnreducedAcc) -> Self;
}
\end{lstlisting}

The \code{UnreducedAcc} associated type is the key to a fast radix-8
matvec: each $8 \times 8$ block accumulates all eight
lifted-twiddle products of an output cell into one length-$31$
unreduced buffer, and folds it back mod $\flift$ exactly once ---
rather than reducing after every one of the eight multiplications.

Implementations exist for \code{DensePolynomial<Int<N>, 16>} (any~$N$,
via crypto-bigint multi-limb arithmetic), the native
\code{DensePolynomial<i64, 16>} and \code{DensePolynomial<i128, 16>},
\code{Gf2\_16} itself (the unlifted reference), and
\code{PackedI64Poly16<BITS>} (which delegates to its inner
\code{DensePolynomial<i64, 16>}). The trait abstraction lets the
\code{ZipTypes} definition site choose the coefficient representation
without touching the kernel.

\subsection{The packed codeword cell}
\label{sec:packed-cw}

A codeword cell is a degree-$16$ polynomial. The shipped \code{Cw}
type is \code{PackedI64Poly16<BITS>}: in memory it is a
\code{DensePolynomial<i64, 16>} (so arithmetic uses native $64$-bit
ALU ops, with headroom for the FFT growth), but its \code{Transcribable}
(wire) encoding writes only \code{BITS} bits per coefficient instead of
the full $64$. The encoder packs $16$ coefficients of \code{BITS} bits
each into $\lceil 16 \cdot \code{BITS} / 8 \rceil$ bytes; the decoder
sign-extends each field back to \code{i64}.

This decouples the in-memory width (fixed at $64$ bits, for safe FFT
arithmetic) from the on-wire width (chosen per use). The shipped Zip$^+$
types use \code{BITS\_CW = 32} for codeword cells --- lossless given the
radix-8 codeword coefficients max near $2^{26}$, with $\sim 6$ bits of
headroom --- and \code{BITS\_CR = 8} for combined-row cells, whose
coefficients are bounded by $\code{batch} \cdot 15$ at $\code{num\_rows}
= 1$.

\subsection{The radix-8 kernel}
\label{sec:kernel}

\code{Radix8FftParams16::new(row\_len, codeword\_len)} asserts
$3 \mid \log_2(\code{codeword\_len})$ and precomputes the
per-meta-stage Vandermonde tables: a
\code{Vec<Vec<[[Gf2\_16; 8]; 8]>>} indexed by
\code{[meta\_stage][block\_idx]}, each entry an $8 \times 8$ matrix
of $\GF(2^{16})$ elements (later lifted at FFT time).

The kernel iterates the $m/3$ meta-stages from the top down; per
meta-stage it processes each block of $8$ in parallel; per block it
runs the $8 \times 8$ matvec using the batched-reduction path:

\begin{lstlisting}[language=Rust]
for j in 0..8 {
    let mut acc = T::fft_unreduced_zero();
    for i in 0..8 {
        let twiddle = matrix[block][j][i];
        if twiddle != Gf2_16::ZERO {
            T::fft_acc_mul_lifted_gf16(&input[i], twiddle, &mut acc);
        }
    }
    output[j] = T::fft_finalize_acc(acc);
}
\end{lstlisting}

\noindent Block-level parallelism (via \code{cfg\_chunks\_mut!} from
\code{zinc-utils}, which compiles to a rayon \code{par\_chunks\_mut}
under the \code{parallel} feature) is safe because each block is a
disjoint slice of the data vector. The zero-twiddle entries are
skipped --- in particular the $X_0 \equiv 1$ column contributes a
plain add.

\subsection{The {\tt LinearCode} impl}
\label{sec:linear-code-impl}

\code{BinaryAddFft16Code<Zt, C, REP, CHECK>} implements
\code{LinearCode<Zt>}:

\begin{lstlisting}[language=Rust]
impl<Zt, C, const REP: usize, const CHECK: bool> LinearCode<Zt>
    for BinaryAddFft16Code<Zt, C, REP, CHECK>
where
    Zt: ZipTypes,
    C: Config16,
    Zt::Cw: FftRingElement16 + FromRef<Zt::Eval>,
    Zt::CombR: FftRingElement16,
{
    const REPETITION_FACTOR: usize = REP;
    fn row_len(&self) -> usize { self.params.row_len }
    fn codeword_len(&self) -> usize { self.params.codeword_len }
    fn encode(&self, row: &[Zt::Eval]) -> Vec<Zt::Cw> {
        let mut data: Vec<Zt::Cw> = row.iter().map(Zt::Cw::from_ref).collect();
        data.resize_with(self.params.codeword_len, Zt::Cw::fft_zero);
        additive_fft_radix8_16(&mut data, &self.params);
        data
    }
    fn encode_wide(&self, row: &[Zt::CombR]) -> Vec<Zt::CombR> {
        let mut data: Vec<Zt::CombR> = row.to_vec();
        data.resize_with(self.params.codeword_len, Zt::CombR::fft_zero);
        additive_fft_radix8_16(&mut data, &self.params);
        data
    }
    fn encode_f<F>(&self, _row: &[F]) -> Vec<F> where ... {
        unimplemented!() // verifier uses encode_wide for proximity
    }
}
\end{lstlisting}

\code{encode} runs the radix-8 FFT on \code{Cw}-typed data;
\code{encode\_wide} runs the identical kernel on \code{CombR}-typed
data for the verifier's proximity check. \code{encode\_f} is part of
the \code{LinearCode} trait but unused by the Zip$^+$ protocol --- the
additive FFT over $\Rlift$ is not naturally expressible over an
arbitrary prime field --- so it is left \code{unimplemented!}.

\subsection{ZipTypes for the binary lane}
\label{sec:zip-types}

The shipped \code{ZipTypes} for this code sets \code{Comb} equal to
\code{CombR}; the prove/verify path then short-circuits the
alpha-sampling logic to $\alpha = [1]$ (because
$\code{Comb::DEGREE\_BOUND} = 0$ when $\code{Comb} = \code{CombR}$,
via a \code{Polynomial<Self>} blanket impl in
\code{poly/src/zero\_degree.rs}). The concrete shape used by the
bench and tests is

\begin{lstlisting}[language=Rust]
impl<...> ZipTypes for BenchAddFft16PolyChalPackedZipTypes<...> {
    const NUM_COLUMN_OPENINGS: usize = NUM_OPENINGS;
    type Eval  = BinaryPoly<16>;
    type Cw    = PackedI64Poly16<BITS_CW>;     // 32 bits/coef on wire
    type Fmod  = Uint<16>;                     // matches F::Modulus width
    type PrimeTest = MillerRabin;
    type Chal  = Bit4Poly16;                   // 4-bit-coef poly challenge
    type Pt    = Bit4Poly16;
    type CombR = PackedI64Poly16<BITS_CR>;     //  8 bits/coef on wire
    type Comb  = Self::CombR;                  // no alpha-projection
    type EvalDotChal     = ScalarProduct;
    type CombDotChal     = ScalarProduct;
    type ArrCombRDotChal = MBSInnerProduct;
}
\end{lstlisting}

\code{Eval = BinaryPoly<16>} is the natural shape of a $1\times$-folded
binary witness. \code{Cw} and \code{CombR} are both
\code{PackedI64Poly16} but with independent on-wire widths, since the
codeword coefficients are large (FFT-amplified) while the combined-row
coefficients stay small ($\le \code{batch}\cdot 15$ at
$\code{num\_rows} = 1$).

\section{Protocol integration}
\label{sec:protocol-integration}

The code is wired end-to-end through Zip$^+$'s commit / prove / verify.
This section describes the verifier-side field and the challenge type
that make that possible. (An earlier draft of this document described
this as an open ``protocol-integration gap''; it is now closed, and the
content below supersedes it.)

\subsection{The obstruction}

The Zip$^+$ prover, at its combined-row / proximity step, computes
inner products of \code{Zt::CombR} slices against verifier-side
field elements. This requires a map
\[
  \mathrm{Zt::CombR} \;\longrightarrow\; F,
\]
i.e.\ $F : \mathrm{FromWithConfig}\langle \&\,\code{Zt::CombR}\rangle$.
For the IPRS / RAA codes \code{CombR} is a scalar \code{Int<N>} and
the map is the standard ``integer mod $p$'' lift into a prime field.
For this code, \code{CombR} is a degree-$16$ polynomial in $\Rlift$,
and no map into a \emph{prime} field $\F_p$ is natural: the encoder's
algebra is polynomial-shaped, so the verifier-side scalar must be too.

\subsection{The verifier-side field $F = \F_p[X]/\flift$}

The fix is to make $F$ itself polynomial-shaped. Define
\[
  F \;=\; \F_p[X]/\flift,
\]
the same quotient shape as $\Rlift$ but with prime-field coefficients
in place of $\Z$-coefficients. The map
\[
  \pi_p : \Rlift \longrightarrow F, \qquad
  \pi_p(\,\textstyle\sum a_i X^i\,) = \sum (a_i \bmod p)\, X^i
\]
is coefficient-wise mod-$p$ reduction. Because $f \equiv \flift
\pmod p$ for the integer-coefficient $\flift$, $\pi_p$ is a
\emph{ring homomorphism} $\Rlift \to F$ --- it commutes with addition,
multiplication, and the $\flift$ reduction --- so it preserves every
algebraic identity the proximity and eval-consistency checks rely on.

When $\flift \bmod p$ is irreducible over $\F_p$, $F$ is a genuine
field, the degree-$16$ extension $\F_{p^{16}}$. The implementation
is \code{Gf2\_16Ext<const P: u64>} in \code{ext\_field.rs}, a
\code{[u64; 16]} coefficient vector mod $p$. It carries
\code{Field::Modulus = Uint<16>} (width-matched to the inner
representation so the transcript's shared modulus / inner buffers
agree) and implements the \code{PrimeField} trait surface that
prove/verify actually exercises --- the trait name is a slight abuse,
but only the \code{Field} operations and config-handling methods are
used; field-specific methods (square roots, $(p-1)/2$) are stubbed.

\paragraph{Prime selection.}
$F$ is a field --- so that \code{Div} / \code{Inv} make sense --- only
when $\flift \bmod p$ is irreducible. \code{ext\_field.rs} ships a
Rabin-style irreducibility test (\code{f\_tilde\_is\_irreducible\_mod\_p})
and a prime scanner. The default prime is
\[
  \code{P\_16\_DEFAULT} \;=\; 2\,147\,482\,921,
\]
the largest prime below $2^{31}$ for which
$\flift = X^{16}+X^5+X^3+X^2+1$ is irreducible mod $p$. A non-ignored
test pins this choice; $P$ is a const-generic, so it can be swapped.

\subsection{The challenge type {\tt Bit4Poly16}}

The row-batching challenge $\code{Chal}$ and evaluation point
$\code{Pt}$ are also polynomial-shaped: \code{Bit4Poly16}, a
$16$-nibble polynomial with $4$-bit integer coefficients in
$\{0,\ldots,15\}$, drawn fresh via Fiat--Shamir. It is backed by a
single \code{i64} ($16 \times 4 = 64$ bits) and has an $8$-byte wire
encoding. The narrow $4$-bit coefficients keep the combined-row
magnitudes small: at $\code{num\_rows} = 1$, a combined-row
coefficient is a sum of $\le \code{batch}$ products of a binary
message bit and a $4$-bit challenge coefficient, hence bounded by
$\code{batch}\cdot 15$ --- which is why \code{BITS\_CR = 8} suffices
for the combined-row wire encoding (\autoref{sec:packed-cw}).

\subsection{The coherence identity}

With $F = \F_p[X]/\flift$ and $\pi_p$ a ring homomorphism, the
protocol's coherence identity holds by algebra: the prover's
combined row, the evaluation, and the tensor coordinates
$(q_0, q_1)$ all live in $F$, and
\[
  \langle \text{combined\_row}, q_1 \rangle_F
  \;=\;
  \langle \text{coeffs}, b \rangle_F
\]
because every quantity is the $\pi_p$-image of the corresponding
$\Rlift$-valued quantity the prover computed, and $\pi_p$ commutes
with the inner product. The transcript serializes the combined-row
vector $b$ as polynomial-valued (packed) entries rather than scalar
field elements.

\subsection{Soundness note}

The proximity argument for the lifted code in $F = \F_p[X]/\flift$ is
not given a formal proof here. The structural ingredients are: the
lifted code reduces mod $2$ to a genuine $\GF(2^{16})$ Reed--Solomon
code (\autoref{sec:lift}), and the Schwartz--Zippel-style argument
the Zip$^+$ proximity test uses has a clean analogue in the extension
$\F_{p^{16}}$ when $\flift \bmod p$ is irreducible. A complete
analysis is left as follow-up (\autoref{sec:future-work}).

\section{Design decisions and performance}
\label{sec:optimization}

This section records the design decisions that shape the shipped
encoder's cost --- both wall-clock and proof size --- and the
reasoning behind each. The single most consequential one is the
radix-8 evaluator; the rest are coefficient-representation and
parallelism choices.

\subsection{Radix-8 evaluator: the coefficient-size win}
\label{sec:opt-radix8}

The dominant design choice is the radix-8 evaluator
(\autoref{sec:radix8}). Its purpose is \emph{not} cache locality ---
it is to control the magnitude of the lifted codeword coefficients.

A codeword cell is $P(v_k) = \sum_i c_i\, X_i(v_k)$ evaluated in
$\Rlift$. How that sum is computed determines how large the
$\Z$-coefficients of the result grow:

\begin{itemize}
  \item \textbf{Radix-2 $\Z$-lift.} The FFT is a chain of $m$
    sequential multiplications in $\Rlift$ by lifted twiddles. Each is
    a $\Z$-arithmetic \code{mul + reduce-mod-}$\flift$; $\Z$-carries
    accumulate across all $m$ stages. Empirically the codeword
    coefficients reach $\sim 2^{47}$ ($49$ signed bits) at
    $\code{codeword\_len} = 4096$ --- they do not fit in \code{i32}
    and need \code{i64} storage and a $\ge 49$-bit wire encoding.
  \item \textbf{Radix-8 with $\GF$-precomputed Vandermonde.} Each
    $8 \times 8$ block matrix is computed once in $\GF(2^{16})$ --- so
    every entry is a $\{0,1\}$-coefficient lifted element --- and the
    FFT composes only $m/3$ such matvecs in $\Rlift$. With a third as
    many carry-accumulating meta-stages, and a $\{0,1\}$-coefficient
    multiplier at each, the codeword coefficients stay near $2^{26}$
    ($28$ signed bits) at the same $\code{codeword\_len} = 4096$.
\end{itemize}

That is a $\sim 21$-bit reduction in coefficient magnitude. It is
\emph{not} free: the radix-8 form does more arithmetic overall
(an $8\times 8$ matvec is $64$ multiplications per block of $8$, versus
$\sim 8$ for three radix-2 stages), so commit time is somewhat higher
than a radix-2 lift would be. The trade is deliberate: the bit-size win
is what lets each codeword cell pack into $32$ bits on the wire
(\autoref{sec:opt-packing}), which is the dominant term of the proof
size. \autoref{sec:measurements} quantifies both sides.

\paragraph{It is a different code.} As noted in \autoref{sec:radix8},
the radix-8-with-$\GF$-lifted-Vandermonde encoder produces a different
lifted code than a radix-2 $\Z$-lift --- the two agree mod $2$ but not
in the higher $\Z$-bits. The small-coefficient variant is the one we
ship.

\subsection{Native {\tt i64} coefficient storage}
\label{sec:opt-native}

Codeword and combined-row cells store their $16$ coefficients as
native \code{i64} (\code{DensePolynomial<i64, 16>}), not
\code{crypto-bigint}'s multi-limb \code{Int<N>}. The radix-8
coefficient bound ($\sim 2^{26}$ at our workload, growing toward
$2^{63}$ only at the largest $\code{codeword\_len} = 2^{16}$) fits
\code{i64}, so the FFT runs on single-instruction $64$-bit ALU adds
instead of generic limb-by-limb carry propagation. The
\code{FftRingElement16} trait keeps the kernel generic, so an
\code{Int<N>} instantiation remains available if a future workload
needs more than $64$ bits per coefficient.

\subsection{Bit-packed wire encoding}
\label{sec:opt-packing}

In-memory width and on-wire width are decoupled. Cells are
\code{i64}-backed in memory (\autoref{sec:opt-native}) but the
\code{Transcribable} encoding of \code{PackedI64Poly16<BITS>}
(\autoref{sec:packed-cw}) writes only \code{BITS} bits per
coefficient.

\begin{itemize}
  \item \textbf{Codeword cells: \code{BITS\_CW = 32}.} The radix-8
    codeword coefficients max near $2^{26}$, so $32$ bits is lossless
    with $\sim 6$ bits of headroom. Column openings are the dominant
    proof component, so this width sets most of the proof size.
  \item \textbf{Combined-row cells: \code{BITS\_CR = 8}.} At
    $\code{num\_rows} = 1$ a combined-row coefficient is bounded by
    $\code{batch}\cdot 15$ (\autoref{sec:protocol-integration}), which
    fits $8$ signed bits at $\code{batch} = 11$. Packing the
    combined row at $8$ bits rather than the full $64$ shrinks that
    proof component $\sim 8\times$.
\end{itemize}

Splitting the two widths matters because the codeword and combined-row
cells have very different magnitude budgets --- codeword coefficients
are FFT-amplified, combined-row coefficients are not.

\subsection{Batched $\flift$-reduction in the matvec}
\label{sec:opt-batched-reduce}

Each $8 \times 8$ block matvec produces one output cell as a sum of
eight lifted-twiddle products. Reducing mod $\flift$ after every one
of the eight multiplications wastes work: the
\code{FftRingElement16::UnreducedAcc} path
(\autoref{sec:fft-ring-element}) accumulates all eight products in a
single length-$31$ unreduced buffer and folds it back mod $\flift$
exactly once per output cell --- one reduction instead of eight.

\subsection{The {\tt parallel} feature}
\label{sec:opt-parallel}

Under the \code{parallel} feature the per-meta-stage block loop is a
rayon \code{par\_chunks\_mut} (via \code{cfg\_chunks\_mut!}). Each
block of $8$ is a disjoint slice, so block-level parallelism is
data-race-free. At commit time there are also \code{batch} independent
row FFTs, giving rayon ample work to distribute. This is the largest
single wall-clock lever; benchmarks in \autoref{sec:measurements} are
all run with \code{parallel} enabled.

\subsection{Zero-twiddle skip}

Vandermonde entries equal to $0 \in \GF(2^{16})$ contribute nothing
and are skipped in the matvec inner loop; the $X_0 \equiv 1$ column,
whose entries are all $1$, contributes a plain add with no
lifted-twiddle multiplication. Small but free.

\section{Measurements}
\label{sec:measurements}

All numbers below come from
\begin{quote}\small
\code{RUSTFLAGS="-C target-cpu=native" cargo bench --bench zip\_plus\_benches}\\
\code{\quad --features parallel,simd,unchecked -- "PolyChal D16"}
\end{quote}
on Apple~M4. The codeword-coefficient measurement is the test
\code{codeword\_coef\_bits} in
\code{zip-plus/src/code/binary\_add\_fft\_16.rs}.

Two configurations are reported, both with $\code{num\_rows} = 1$ and
$\code{batch} = 11$:
\begin{itemize}
  \item \textbf{D16 R4}: $\code{REP} = 4$, $\code{num\_vars} = 10$
    ($\code{poly\_size} = 2^{10}$, $\code{codeword\_len} = 4096$).
  \item \textbf{D16 R8}: $\code{REP} = 8$, $\code{num\_vars} = 9$
    ($\code{poly\_size} = 2^{9}$, $\code{codeword\_len} = 4096$).
\end{itemize}
The radix-8 evaluator requires $3 \mid \log_2(\code{codeword\_len})$,
which both configurations satisfy ($\log_2 4096 = 12$). The D16 R4
configuration commits the same total bit-count
($2^{10}\cdot 16 = 2^{14}$) as the protocol's $1\times$-folded binary
lane at $\code{protocol\_num\_vars} = 9$.

\subsection{Codeword coefficient magnitude}
\label{sec:meas-coef}

The defining measurement for this code. On random
\code{BinaryPoly<16>} input at $\code{codeword\_len} = 4096$:

\begin{center}\small
\begin{tabular}{lrr}
  \toprule
  Evaluator & Max $|$coef$|$ & Bits incl.\ sign \\
  \midrule
  Radix-8 (shipped)           & $\sim 2^{26.4}$ & $28$ \\
  Radix-2 $\Z$-lift (historic)& $\sim 2^{47.1}$ & $49$ \\
  \bottomrule
\end{tabular}
\end{center}

\noindent The radix-8 evaluator's $\sim 21$-bit reduction
(\autoref{sec:opt-radix8}) is what makes a $32$-bit on-wire codeword
cell lossless. A radix-2 $\Z$-lift would need a $\ge 49$-bit encoding,
inflating the dominant proof component by $\sim 1.5\times$.

\subsection{Commit / prove / verify timing}

Representative wall-clock figures (Apple~M4, \code{parallel} +
\code{simd} + \code{unchecked}; the M-series exhibits noticeable
run-to-run variance, so treat these as $\pm 30\%$):

\begin{center}\small
\begin{tabular}{lrrr}
  \toprule
  Configuration & Commit & Prove & Verify \\
  \midrule
  D16 R4 & $\sim 3$--$7$\,ms & $\sim 1.8$\,ms & $\sim 1.1$--$1.7$\,ms \\
  D16 R8 & $\sim 7$\,ms       & $\sim 1.2$\,ms & $\sim 1.5$\,ms \\
  \bottomrule
\end{tabular}
\end{center}

\noindent Commit is the heaviest step --- it runs the radix-8 FFT on
all \code{batch} rows --- and is where the radix-8 form costs more raw
arithmetic than a radix-2 lift would (\autoref{sec:opt-radix8}). Prove
and verify are dominated by the transcript / Merkle-opening work and
are largely independent of the FFT.

\subsection{Proof size}

Total proof size (commitment + transcript), raw and after zstd-3
compression:

\begin{center}\small
\begin{tabular}{lrr}
  \toprule
  Configuration & Raw & zstd-3 \\
  \midrule
  D16 R4 & $\sim 176$\,KiB & $\sim 127$--$138$\,KiB \\
  D16 R8 & $\sim 117$\,KiB & $\sim 83$\,KiB \\
  \bottomrule
\end{tabular}
\end{center}

\noindent The proof is dominated by column openings: each of the
\code{NUM\_COLUMN\_OPENINGS} openings reveals $\code{batch}$ codeword
cells, each a $16$-coefficient polynomial packed at \code{BITS\_CW}$=32$
bits per coefficient. Rate $1/8$ (D16 R8) uses fewer openings ($100$
versus $147$) and so produces a smaller proof, at the cost of a longer
codeword and heavier commit.

For reference, the IPRS code at the same SHA-256 binary-lane
dimensions ($\code{num\_vars} = 9$, $\code{batch} = 11$, run as the
\code{Zip+ SHA-style (IPRS)} bench group) commits the same
$2^{14}$-bit workload and produces a $\sim 479$\,KiB raw /
$\sim 288$\,KiB zstd proof. The D16 R4 polychal code is $\sim 2.7\times$
smaller raw and $\sim 2.2\times$ smaller compressed at that workload ---
the binary-additive-FFT structure plus the radix-8 small-coefficient
codeword is a substantially tighter fit for binary witnesses than
running them through the $\F_{65537}$ IPRS encoder.

\subsection{Why the lifted codeword is compressible}

A codeword cell holds $16$ \code{i64} coefficients, but the radix-8
coefficients occupy only $\sim 28$ bits each, and \code{BITS\_CW}$=32$
already trims the on-wire encoding to $32$ of the $64$ in-memory bits.
Beyond that, many evaluation points carry significant sparsity (the
input is zero-padded from $\code{row\_len}$ to
$\code{codeword\_len} = \code{REP}\cdot\code{row\_len}$), and the
$\{0,1\}$-coefficient Vandermonde structure leaves predictable
low-entropy patterns. zstd-3 recovers a further $\sim 25$--$30\%$ on
top of the bit-packing.

\subsection{Test coverage}

The following tests in
\code{zip-plus/src/code/binary\_add\_fft\_16/} exercise the
implementation:

\begin{itemize}
  \item \code{basis::tests::*}: \code{Gf2\_16} arithmetic
    identities, the Cantor recurrence $\bb_i^2 + \bb_i = \bb_{i-1}$,
    basis linear independence, the half-trace solver.
  \item \code{params::tests::*}: subspace polynomial monicity,
    $s_i(\bb_i) = 1$, $s_i$ vanishes on $V_i$, the Cantor recurrence
    at evaluation points.
  \item \code{ring\_ops::tests::*}: $\flift$-reduction correctness,
    the $\Rlift \to \GF(2^{16})$ commutative diagram for
    multiplication, the lifted-twiddle multiplication contract.
  \item \code{radix8::tests::*}: the radix-8 parameter constructor
    accepting / rejecting $\log_2(\code{codeword\_len})$ by
    divisibility-by-$3$.
  \item \code{ext\_field::tests::*}: $\GF(2^{16})$-extension
    arithmetic, and \code{default\_prime\_makes\_f\_tilde\_irreducible}
    pinning \code{P\_16\_DEFAULT}.
  \item \code{linear\_code\_tests::*}: the \code{LinearCode<Zt>}
    encode smoke test; the full commit / prove / verify roundtrip
    (\code{commit\_prove\_verify\_polychal\_smoke\_16}); tamper
    detection (\code{polychal\_tamper\_detection\_16}); and the
    codeword-coefficient measurement (\code{codeword\_coef\_bits}).
\end{itemize}

All tests pass under \code{cargo test -p zip-plus --features
parallel,simd,unchecked}.

\section{Open questions and future work}
\label{sec:future-work}

\subsection{Soundness analysis}

The implementation is complete and exercised end-to-end, but the
soundness of the lifted code in the verifier-side field
$F = \F_p[X]/\flift$ is not formally proved here. The structural
ingredients are in place --- the lifted code reduces mod $2$ to a
genuine $\GF(2^{16})$ Reed--Solomon code (\autoref{sec:lift}), which
is MDS with distance $N - n + 1$ --- but a complete analysis should:
\begin{itemize}
  \item Establish that the lifted code's effective minimum distance,
    as seen by the proximity test, matches (or closely tracks) the
    underlying $\GF(2^{16})$ RS distance, accounting for the
    structural lift's torsion.
  \item Carry the Zip$^+$ proximity argument through the
    $F = \F_p[X]/\flift$ setting. The existing argument uses a
    Schwartz--Zippel bound over a prime field; the analogue over the
    degree-$16$ extension $\F_{p^{16}}$ (valid when $\flift \bmod p$
    is irreducible, as it is for \code{P\_16\_DEFAULT}) is the natural
    target.
  \item Bound the soundness-error contribution of the radix-8
    evaluator specifically. The radix-8-with-$\GF$-lifted-Vandermonde
    code differs from a radix-2 $\Z$-lift in its higher $\Z$-bits
    (\autoref{sec:radix8}); the analysis should confirm the proximity
    argument is insensitive to that difference, since both reduce mod
    $2$ to the same $\GF(2^{16})$ code.
\end{itemize}

\subsection{Rate / opening-count tuning}

The shipped configurations use $\code{REP} = 4$ ($147$ column
openings) and $\code{REP} = 8$ ($100$ openings), matching the
protocol's IPRS opening table. The radix-8 evaluator constrains
$\log_2(\code{codeword\_len})$ to be a multiple of $3$, which couples
the admissible $(\code{row\_len}, \code{REP})$ pairs; e.g.\ a
bit-budget-matched $\code{REP} = 8$ run at $\code{num\_vars} = 10$
would need $\log_2(\code{codeword\_len}) = 13$, which the radix-8
constructor rejects. A radix-$2$/radix-$8$ \emph{mixed} evaluator
(one or two radix-2 stages plus radix-8 meta-stages) would lift this
divisibility constraint, at the cost of a slightly more intricate
kernel and a small reintroduction of radix-2 carry growth.

\subsection{Tighter codeword packing}

\code{BITS\_CW = 32} carries $\sim 6$ bits of headroom over the
empirical $\sim 2^{26}$ codeword-coefficient maximum at
$\code{codeword\_len} = 4096$. A workload-specific bound could push
\code{BITS\_CW} down toward $28$, shrinking the dominant proof
component a further $\sim 10\%$ --- at the cost of a non-power-of-two
field width and a tighter (workload-dependent) safety margin. The
coefficient magnitude grows with $\code{codeword\_len}$ (it approaches
$2^{63}$ at the maximal $\code{codeword\_len} = 2^{16}$, see
\autoref{sec:meas-coef}), so any tightened \code{BITS\_CW} must be
chosen per codeword length.

\subsection{Performance follow-ups}

\begin{itemize}
  \item \textbf{Within-FFT parallelism for small batches.} At
    $\code{batch} = 1$ the row-level rayon parallelism vanishes and
    only the per-meta-stage block loop remains; rayon's adaptive
    splitter may under-extract parallelism for small block counts.
  \item \textbf{SIMD matvec.} The $8 \times 8$ block matvec multiplies
    by $\{0,1\}$-coefficient lifted twiddles --- shift-and-add over
    \code{i64} coefficients --- which is a natural target for an
    explicit SIMD inner loop on hosts with wide vector registers.
  \item \textbf{Commit-side cost of radix-8.} The radix-8 form does
    more raw arithmetic than a radix-2 lift (\autoref{sec:opt-radix8}).
    The current trade favours proof size; a deployment that weights
    commit time more heavily could expose the radix as a tunable knob.
\end{itemize}

\subsection{Larger / other $\GF(2^k)$}

The implementation is fixed at $k = 16$, matching the protocol's
$1\times$-folded binary lane (\code{BinaryPoly<16>}). The algebra is
identical for other even $k$; only the precomputation changes ---
the irreducible $f$, the Cantor basis size, and the
\code{reduce\_mod\_ftilde} fold pattern. A $\GF(2^{32})$ instantiation
would serve the \emph{unfolded} binary lane directly (no protocol-side
folding), at the cost of larger codeword coefficients and a
$\GF(2^{32})$ evaluation domain.

\subsection{Code-survey reference}

For readers exploring the source, the key files are:

\begin{itemize}\small
  \item \code{zip-plus/src/code/binary\_add\_fft\_16.rs}: top-level,
    \code{BinaryAddFft16Code}, \code{LinearCode<Zt>} impl, smoke /
    roundtrip / measurement tests.
  \item \code{.../binary\_add\_fft\_16/basis.rs}:
    \code{Gf2\_16}, Cantor basis.
  \item \code{.../binary\_add\_fft\_16/params.rs}:
    \code{Config16}, precomputed basis / subspace polynomials.
  \item \code{.../binary\_add\_fft\_16/ring\_ops.rs}:
    \code{FftRingElement16}, lifted-ring arithmetic.
  \item \code{.../binary\_add\_fft\_16/radix8.rs}:
    \code{Radix8FftParams16}, the radix-8 kernel.
  \item \code{.../binary\_add\_fft\_16/ext\_field.rs}:
    \code{Gf2\_16Ext<P>} ($F = \F_p[X]/\flift$), \code{Bit4Poly16},
    \code{P\_16\_DEFAULT}, the irreducibility test.
  \item \code{.../binary\_add\_fft\_16/packed\_cw.rs}:
    \code{PackedI64Poly16<BITS>}, the bit-packed wire encoding.
  \item \code{zip-plus/benches/zip\_plus\_benches.rs}: the bench
    harness --- RAA / IPRS references plus the D16 R4 / R8 polychal
    groups.
\end{itemize}


\end{document}
